\documentclass[11pt]{article}

\usepackage[margin=1in]{geometry}
\usepackage{booktabs}
\usepackage{array}
\usepackage{amsmath}
\usepackage{hyperref}
\hypersetup{colorlinks=true, linkcolor=black, urlcolor=blue}

\newcommand{\code}[1]{\texttt{\small #1}}

\title{\texttt{aivreg} 1.1.2: Package Update Report}
\date{31 August 2026}

\begin{document}
\maketitle

\section{What was merged}

I merged Alex's working candidate, \code{aivreg\_working\_candidate\_1\_1\_1\_20260825.ado},
with Josh's most recent updates on branch \code{gmm\_fix}, and called the result
\code{aivreg\_merged\_draft\_1\_1\_2\_20260830.ado} (version 1.1.2).
Table~\ref{tab:files} lists the versions used throughout this note.

\begin{table}[h]
\centering
\caption{Versions compared}
\label{tab:files}
\begin{tabular}{lll}
\toprule
Label & File & Role \\
\midrule
Baseline  & \code{aivreg.ado} (1.0.0)                        & Repository state before the merge \\
Candidate & \code{aivreg\_working\_candidate\_1\_1\_1\_20260825.ado} & Alex's reviewed candidate \\
Merged    & \code{aivreg\_merged\_draft\_1\_1\_2\_20260830.ado}     & 1.1.1 $+$ Josh's updates \\
\bottomrule
\end{tabular}
\end{table}


\section{Weighted and clustered GMM}
\label{sec:weights}

Let $w_i$ be the probability weight, $W=\sum_i w_i$, $N$ the number of observations,
$\hat{w}_i = w_i N / W$ the weight normalised to mean one, and $\varepsilon_i$ the
observation's moment contribution. Table~\ref{tab:weights} states what each version does.

\begin{table}[h]
\centering
\caption{Weight and cluster handling in the GMM engine}
\label{tab:weights}
\renewcommand{\arraystretch}{1.25}
\begin{tabular}{p{0.29\textwidth} p{0.29\textwidth} p{0.29\textwidth}}
\toprule
 & Baseline & Candidate / merged \\
\midrule
Weight in the point estimate & $w_i/W$ (linear) & $w_i/W$ (linear; unchanged) \\
Weight in the covariance ``meat'' & $w_i/W$ --- linear & $\hat{w}_i^{\,2}$ --- squared \\
Finite-sample factor, unclustered & $W/(W-K)$ & $N/(N-K)$ \\
Finite-sample factor, clustered & $\dfrac{G}{G-1}\cdot\dfrac{W-1}{W-K}$ & $\dfrac{G}{G-1}\cdot\dfrac{N-1}{N-K}$ \\
Cluster score sum & $\sum_{i \in g}\sqrt{w_i/W}\,\varepsilon_i$ & $\sum_{i \in g}\hat{w}_i\,\varepsilon_i$ \\
Clustered $\bar{g}$ & cluster weight share $\times$ cluster sums (weights counted twice) & $N^{-1}\sum_g\sum_{i \in g}\hat{w}_i\varepsilon_i$, consistent with the unclustered $\bar{g}$ \\
Rescaling all weights by a constant & changes the standard errors and $J$ & leaves $\hat{\beta}$, $V$ and $J$ unchanged \\
Zero or negative weights & silently dropped & error 459 \\
\code{[aw=]} or \code{[fw=]} on \code{gmm}/\code{2sls} & accepted, treated as \code{pw} & error 101 \\
\bottomrule
\end{tabular}
\end{table}

\noindent In words: the baseline folded $\sqrt{w_i/W}$ into each score and then formed the
outer product, so the weight entered the covariance linearly rather than squared; and it
used the sum of weights in place of the observation count in the finite-sample factors,
so multiplying every weight by a constant changed the reported standard errors. On the
clustered path the weight was applied once inside the cluster sum and again as a cluster
weight share when forming $\bar{g}$. The candidate applies mean-one normalised weights
linearly to the moment and in squared form to the covariance, bases the finite-sample
factors on $N$ and $G$ only, and builds the clustered $\bar{g}$ directly from this call's
own cluster sums.

\noindent The effect is visible in tests \code{A11} and \code{A12} of
Table~\ref{tab:results}, which are the same specification with every weight multiplied by
$17$. The baseline standard errors differ in the fifth significant digit
($0.0267815$ against $0.0267743$); the merged standard errors agree to roughly twelve
digits.

\section{Degrees of freedom for the Hansen $J$ test}

With $k$ instruments (amenities and controls) and $L$ anti-IVs, the moment system has
$m = L(k+2)$ conditions and $p = k + 2L$ parameters. \textbf{The adopted rule is the
conventional one,}
\[
  \mathrm{df} \;=\; m - p \;=\; L(k+2) - (k+2L) \;=\; k(L-1),
\]
with $J = N\,\bar{g}'W_2\,\bar{g}$ evaluated at the second-step estimate and the
efficient weighting matrix $W_2$, and $p$-value $= \Pr[\chi^2_{\mathrm{df}} > J]$.
An earlier proposal set the degrees of freedom from numerical matrix ranks,
$\operatorname{rank}(W_2) - \operatorname{rank}(X'W_2X)$. That rule is \emph{not} used.

\noindent Ranks are retained only as a validity gate. Before reporting $J$, the command
checks $\operatorname{rank}(W_2) = m$ and $\operatorname{rank}(X'W_2X) = p$. If either
check fails, $J$ is suppressed entirely, a warning is printed, and \code{e(J\_status)}
records the suppression; no generalised rank-based degrees of freedom are ever
substituted. This is the conservative reading: a rank-deficient moment system has no
valid $\chi^2$ reference, so the statistic is withheld rather than reported against an
adjusted distribution.

\noindent The $J$ statistic is reported only on the efficient two-step GMM path. One-step
GMM, \code{2sls}, and just-identified problems print an explicit note in place of a
statistic, rather than a criterion value that would be read as a Hansen test.

\section{Regression test results}

The suite in \code{tests/aivreg\_test\_suite.do} runs 29 tests against each staged version
and records the return code, the first coefficient, its standard error, and the $J$
statistic with its degrees of freedom. Section~A covers ordinary documented usage;
Section~B covers edge cases and the defects under review. Table~\ref{tab:tests} says
what question each test answers; Table~\ref{tab:results} reports both runs.

\begin{table}[htbp]
\centering
\caption{What each test asks}
\label{tab:tests}
\footnotesize
\begin{tabular}{ll}
\toprule
Test & Question it answers \\
\midrule
\multicolumn{2}{l}{\textit{Section A --- ordinary documented usage}} \\
\code{A1}   & Does the single-anti-IV ratio path work with the default Anderson--Rubin SEs? \\
\code{A2}   & Does the ratio path work with asymptotic SEs? \\
\code{A3}   & Does the ratio path work with bootstrap SEs (seeded)? \\
\code{A4}   & Do factor-variable controls with \code{if} work on the ratio path? \\
\code{A5}   & Do fixed effects work on the ratio path? \\
\code{A6}   & Does a multi-anti-IV call switch to GMM automatically? \\
\code{A7}   & Does two-step GMM run with multiple anti-IVs? \\
\code{A8}   & Does one-step GMM run? \\
\code{A9}   & Does 2SLS run? \\
\code{A10}  & Does clustering (on \code{fips}) work under GMM? \\
\code{A11}  & Do probability weights work under GMM? \\
\code{A12}  & Same call as \code{A11} with all weights $\times 17$: are results scale invariant? \\
\code{A13}  & Does the ratio path still accept its documented \code{[aw=]} syntax? \\
\midrule
\multicolumn{2}{l}{\textit{Section B --- edge cases and known defects}} \\
\code{B1}   & Is an \code{[in]} range honoured on the ratio path? \\
\code{B2}   & Is an \code{[in]} range honoured on the GMM path? \\
\code{B3}   & Are \code{if} and \code{[in]} combined honoured on the GMM path? \\
\code{B4}   & Is an explicit \code{[aw=]} rejected loudly on the GMM path? \\
\code{B5}   & Is an explicit \code{[fw=]} rejected loudly on the GMM path? \\
\code{B6}   & Are zero or negative weights refused loudly? \\
\code{B7}   & Are observations with missing weights dropped from the sample? \\
\code{B8}   & Does fewer than two clusters produce a clean error? \\
\code{B9}   & Do 32-character variable names work on the ratio path? \\
\code{B10}  & Do 32-character variable names work on the GMM path? \\
\code{B11}  & Does a failed call avoid leaving the previous call's results in \code{e()}? \\
\code{B12}  & Do fixed effects, clustering and weights work together? \\
\code{B13}  & Is $J$ suppressed on a just-identified problem (single anti-IV via \code{gmm})? \\
\code{B14}  & Does a collinear anti-IV fail rather than silently collapse the moment system? \\
\code{B15}  & Does a mistyped estimator (\code{gmmm}) error rather than run? \\
\code{B16}  & Do non-consecutive integer fixed-effect levels work? \\
\bottomrule
\end{tabular}
\end{table}

% Notes typeset inside the generated results float (see make_table.py).
\newcommand{\resultsnotes}{\textit{Notes.} ``--'' means the quantity was not returned.
Baseline rows with a non-zero return code still report a coefficient because the
baseline left the previous call's results in \code{e()}; the merged draft clears them.
Return code 101 is the refusal of \code{[aw=]} or \code{[fw=]} on the GMM path, 459 the
refusal of non-positive weights, 430 the singular anti-IV system, and 111 and 198 on the
baseline the long-name failures the merge repairs; 198 on merged \code{B15} is the
intended rejection of the mistyped estimator.}

\begin{table}[htbp]
\centering
\caption{Baseline versus merged candidate, 29 regression tests}
\label{tab:results}
\footnotesize
\setlength{\tabcolsep}{4pt}
\begin{tabular}{l cc cc cc cc}
\toprule
 & \multicolumn{2}{c}{Return code} & \multicolumn{2}{c}{Coefficient}
 & \multicolumn{2}{c}{Std.\ error} & \multicolumn{2}{c}{$J$ (df)} \\
\cmidrule(lr){2-3}\cmidrule(lr){4-5}\cmidrule(lr){6-7}\cmidrule(lr){8-9}
Test & Base & Merg & Base & Merg & Base & Merg & Base & Merg \\
\midrule
A1\_ratio\_default & 0 & 0 & -1.1451 & -1.1451 & 0.1011 & 0.1011 & -- & -- \\
A2\_ratio\_asymp & 0 & 0 & -1.1451 & -1.1451 & 0.1091 & 0.1091 & -- & -- \\
A3\_ratio\_boot & 0 & 0 & -1.1451 & -1.1451 & 0.1154 & 0.1154 & -- & -- \\
A4\_ratio\_controls\_if & 0 & 0 & -0.5250 & -0.5250 & 0.0204 & 0.0204 & -- & -- \\
A5\_ratio\_fe & 0 & 0 & -0.5250 & -0.5250 & 0.0204 & 0.0204 & -- & -- \\
A6\_multiamenity\_gmm & 0 & 0 & -0.5355 & -0.5355 & 0.0213 & 0.0213 & -- & -- \\
A7\_gmm\_twostep & 0 & 0 & -0.5369 & -0.5369 & 0.0275 & 0.0275 & $1.86\times 10^{-5}$ (--) & 0.3987 (2) \\
A8\_gmm\_onestep & 0 & 0 & -0.4237 & -0.4237 & 0.1856 & 0.1856 & 0.0074 (--) & -- \\
A9\_2sls & 0 & 0 & -0.5385 & -0.5385 & 0.0279 & 0.0279 & 0.9993 (--) & -- \\
A10\_gmm\_cluster & 0 & 0 & -0.5222 & -0.5222 & 0.0590 & 0.0590 & $6.70\times 10^{-7}$ (--) & 8.7252 (2) \\
A11\_gmm\_pweight\_bare & 0 & 0 & -0.5320 & -0.5319 & 0.0268 & 0.0278 & $2.16\times 10^{-5}$ (--) & 0.4773 (2) \\
A12\_gmm\_pweight\_x17 & 0 & 0 & -0.5320 & -0.5319 & 0.0268 & 0.0278 & $2.16\times 10^{-5}$ (--) & 0.4773 (2) \\
A13\_ratio\_aw\_bracket & 0 & 0 & -0.8091 & -0.8091 & 0.0345 & 0.0345 & -- & -- \\
B1\_ratio\_in & 0 & 0 & -2.9689 & -2.9689 & 0.1485 & 0.1485 & -- & -- \\
B2\_gmm\_in & 0 & 0 & -1.6825 & -1.6825 & 0.1161 & 0.1161 & 0.0082 (--) & 339.3547 (2) \\
B3\_gmm\_if\_in & 0 & 0 & -1.6289 & -1.6289 & 0.1163 & 0.1163 & 0.0008 (--) & 9.8732 (2) \\
B4\_gmm\_aw\_bracket & 101 & 101 & -1.6289 & -- & 0.1163 & -- & 0.0008 (--) & -- \\
B5\_gmm\_fw\_bracket & 101 & 101 & -1.6289 & -- & 0.1163 & -- & 0.0008 (--) & -- \\
B6\_gmm\_negzero\_weight & 0 & 459 & -0.5047 & -- & 0.0267 & -- & $2.54\times 10^{-5}$ (--) & -- \\
B7\_gmm\_missing\_weight & 0 & 0 & -0.5527 & -0.5527 & 0.0295 & 0.0306 & $2.30\times 10^{-6}$ (--) & 0.0270 (2) \\
B8\_gmm\_two\_clusters & 430 & 430 & 0.0000 & 0.0000 & 0.0000 & 0.0000 & 0.2499 (--) & -- \\
B9\_ratio\_longnames & 111 & 0 & 3.0944 & -0.8396 & 0.1101 & 0.0392 & -- & -- \\
B10\_gmm\_longnames & 198 & 0 & -0.8386 & -0.8394 & 0.0687 & 0.0460 & $9.01\times 10^{-8}$ (--) & 0.0002 (1) \\
B11\_state\_after\_error & 0 & 459 & -0.5047 & -- & 0.0267 & -- & $2.54\times 10^{-5}$ (--) & -- \\
B12\_gmm\_fe\_clus\_w & 0 & 0 & -0.5279 & -0.5279 & 0.0579 & 0.0586 & $1.09\times 10^{-8}$ (--) & $2.66\times 10^{-6}$ (1) \\
B13\_gmm\_justidentified & 0 & 0 & -0.5293 & -0.5293 & 0.0279 & 0.0279 & -- & -- \\
B14\_gmm\_collinear\_aiv & 430 & 430 & 0.0356 & 0.0356 & 0.0053 & 0.0053 & 0.9994 (--) & -- \\
B15\_bad\_estimator & 0 & 198 & -0.8386 & -- & 0.0687 & -- & $9.01\times 10^{-8}$ (--) & -- \\
B16\_fe\_nonconsecutive & 0 & 0 & -0.5250 & -0.5250 & 0.0204 & 0.0204 & -- & -- \\
\bottomrule
\end{tabular}

\vspace{0.6em}
\begin{minipage}{\textwidth}
\footnotesize
\resultsnotes % defined in main.tex just before the \input
\end{minipage}
\end{table}


\subsection*{Reading the table}

\begin{itemize}
  \item \textbf{Hansen $J$.} The $J$ calculation is now correct. The baseline returned a
        degenerate near-zero $J$ with no degrees of freedom on every over-identified
        call; the merged draft reports a proper statistic with its degrees of freedom
        (\code{A7}, \code{A10}) and suppresses $J$ where the test is not defined.
  \item \textbf{Loud failures.} \code{B4}, \code{B5}, \code{B6} and \code{B11} now exit
        with an error and leave no results, where the baseline either ran or reported
        stale values.
  \item \textbf{Long names.} \code{B9} and \code{B10} fail under the baseline
        (return codes 111 and 198) and succeed under the merge.
\end{itemize}

\end{document}
